\tightsection{How Predictable Are Video Metrics?}
\label{predictability}

Before we delve into algorithms, we are going to answer the question of how preditable are the video metrics using the session-level attributes. Intuitively, if all sessions with the same attribute values exhibit similar video quality, then the set of attributes are predictive, and vice versa. In this section, we will present an empirical predictability analysis based on real-world data. Such analysis also establishes an upper bound that can help us evaluate the prediction algorithm presented in \Section~\ref{subsec:upperbound}.

\tightsubsection{Definition and model}
\label{subsec:upperbound}

We first define predictability. We start with defining the {\it prediction error}.  Given a session
(which we will call the {\it session under prediction}), let $q$ be
the actual quality and $p$ be the predicted quality. Then, the
prediction error $e_{p,q}$ is defined as
\begin{packeditemize}
	\item $e_{p,q}=|p-q|$, in case of buffering ratio and start failure (1 is failure and 0 is success);
	\item $e_{p,q}=|p-q|/q$, in case of average bitrate and join time.
\end{packeditemize}

For a set of sessions under prediction $S=\{s_1,\dots,s_n\}$ let $P=\{p_i\}$ and $Q=\{q_i\}$ be their predicted and actual qualities. respectively ($p_i$ and $q_i$ are the predicted and actual quality of $s_i$). The overall prediction error is then $E_{P,Q}=\left(\frac{1}{n}\sum e_{p_i,q_i}^2\right)^{1/2}$.

% xil: I feel like this is a simple concept that can be captured with plain English.
%We define {\it attribute combination} ({\it AC}) $g$ to be a set of attributes. Given an AC $g$, function $v_g$ takes a session $s$ as input and returns an array of values of $s$ on each attribute in $g$. For instance, if $g=[\textrm{ASN, CDN}]$, $v_g(s)=[ASN_s,CDN_s]$ where the client of $s$ belongs to $ASN_s$ and the server belongs to $CDN_s$. An {\it attribute-based prediction algorithm} $P$ \ion{$P$ is also the predicted quality vector, so need to change the notation.} is defined by an AC $g$. $P$ takes a session $s$ as input, and returns $P_g(s)$ as the predicted quality. Note that such an algorithm will provide the same prediction to all sessions with identical attribute values. This means that $P$ cannot provide perfect prediction unless we identify all possible attributes that may impact a sessions's quality, which in practice is infeasible.

We define {\it attribute combination} ({\it AC}) $g$ to be a set of attributes, e.g. $[ASN, CDN]$. An {\it attribute-based prediction algorithm} takes a session $s$ and an AC $g$ as input, and returns the predicted quality. 

%\xil{the point of the next sentence is unclear.} Note that such an algorithm will provide the same prediction to all sessions with identical attribute values. This means that $P$ cannot provide perfect prediction unless we identify all possible attributes that may impact a sessions's quality, which in practice is infeasible.

The {\it predictability} of a particular AC over a set of sessions is the minimal overal prediction error that can be achieved, i.e., $\textrm{argmax}E_{P,Q}=\textrm{argmax}_p\left(\frac{1}{n}\sum e_{p,q_i}^2\right)^{1/2}$
For example, the optimal buffering ratio prediction for a session from a group is to predict using the mean of buffering ratio from that group. The predictability quantifies the dispersion in quality of the sessions that an AC cannot differentiate. Ideally, if the attributes in AC selected reflect all the factors that determine the quality of a session, then the sessions in an identical group of $g$ should produce the same quality and the upper bound is exactly one. 

%Therefore, the predicability of $g$ over a session set $S$ is the minimal overall prediction error $R_g(S)=E_{P^*,Q}$ where $P^*$ gives the same optimal prediction $p^*$ for each session in $S$. This definition can be extended to any set of sessions $S$ by dividing sessions into identical groups of $g$, and use the optimal prediction for each identical group to be the prediction of sessions in it, i.e., $R_g(S)=E_{P^*,Q}$ where prediction of $s_i$ in $P^*$ is the optimal prediction of the identical group that $s_i$ belongs to.\ion{This is a very confusing paragraph.}


%The {\it predictability} of AC $g$ over a set of sessions $S$ is the minimal overal prediction error over all prediction algorithms on $g$, i.e., $p^*=\textrm{argmax}E_{P,Q}=\textrm{argmax}_p\left(\frac{1}{n}\sum e_{p,q_i}^2\right)^{1/2}$. For example, the optimal buffering ratio prediction for a session from a group is to predict using the mean of buffering ratio from that group. Therefore, the predicability of $g$ over a session set $S$ is the minimal overall prediction error $R_g(S)=E_{P^*,Q}$ where $P^*$ gives the same optimal prediction $p^*$ for each session in $S$. This definition can be extended to any set of sessions $S$ by dividing sessions into identical groups of $g$, and use the optimal prediction for each identical group to be the prediction of sessions in it, i.e., $R_g(S)=E_{P^*,Q}$ where prediction of $s_i$ in $P^*$ is the optimal prediction of the identical group that $s_i$ belongs to.\ion{This is a very confusing paragraph.} 

%The predictability quantifies the dispersion in quality of the sessions that an AC cannot differentiate. Ideally, if the attributes in AC $g$ selected reflect all the factors that determine the quality of a session, then the sessions in an identical group of $g$ should produce the same quality and the upper bound is exactly one. 

\jc{The current definition predictability is weird as it gives smaller value to more predictable set of sessions. }


\begin{figure*}[t!]
\centering
\subfigure[Predictability of temporal attributes]
{
	\includegraphics[width=0.2\textwidth]{figures/placeholder.pdf}
	\label{subfig:predictability:temporal}
}
\subfigure[Predictability of spatial attributes]
{
	\includegraphics[width=0.2\textwidth]{figures/placeholder.pdf}
	\label{subfig:predictability:temporal}
}
\subfigure[Rank of predictability among ASN]
{
	\includegraphics[width=0.2\textwidth]{figures/placeholder.pdf}
	\label{subfig:predictability:asn}
}
\subfigure[Rank of predictability among Site]
{
	\includegraphics[width=0.2\textwidth]{figures/placeholder.pdf}
	\label{subfig:predictability:site}
}
\tightcaption{Quantifying predictability over our dataset.}
\label{fig:predictability}
\end{figure*}


\tightsubsection{Empirical predictability analysis}

\xil{need some clean up after figures are in, describe data set, set up ``time is an attribute" (maybe even before)}

In this section, we quantify the predictability of the session attributes in our dataset. Such analysis provides an upper-bound for any prediction algorithm, and it may not be a tight upper-bound. Despite these limitations, predictability is a valuable bound to compare prediction algorithms agianst. We consider the set of spatial attributes introduced in \Section~\ref{subsec:dataset} as well as a set of time intervals as temporal attributes.

\myparatight{Temporal attributes} Figure~\ref{subfig:predictability:temporal} shows the predictability of different temporal attributes (time granular) combined with all spatial attributes. An example of an AC with a $5$-minute temporal attribute is [5-minute, ASN, Object, Site, Initial CDN, Initial Bitrate, OS, ConnectionType]. Each point in the figure represents the overall predictability for a given time period. The predictability tends to be better for small intervals. As the interval length increases the predictability stabilizes.
%The figure shows that with different granulars do have impact on predictability. Smaller intervals are more predictable, but the predictability becomes similar after granular is less than about 5 minutes.

\myparatight{Spatial attributes} Figure~\ref{subfig:predictability:temporal} shows the predictability of different spatial attributes combined with the $5$ minute temporal attribute. Again, the finest spatial attribute combination is [ASN, Object, Site, Initial CDN, Initial Bitrate, OS and ConnectionType].  Instead of showing all ACs with these attributes (in total $2^7-1=127$ ACs), we group them by their number of attributes and show the ones with best predictability. This illustrates the impact of the number of attributes on predictability. The results show clear improvements as we use more attributes, though we reach diminishing return when we approach seven attributes.

\myparatight{Different partitions} The above results show that predictability varies significantly with the attributes we consider. Next, we present the predictability also varies across different partitions of the same AC (e.g., ASN, Site). Figure~\ref{subfig:predictability:asn} (/\ref{subfig:predictability:site}) shows the distribution of the mean predictability among top-$k$ and last-$k$ ASNs (/Sites). These results confirm that predictability vary significantly across different partitions. 
%More importantly, this suggests that for some classes of clients, the video quality is more predictable than for other classes of clients by using the attributes we collect.

\comment{
\tightsubsection{Summary of observations}
\begin{packedenumerate}
	\item The overall prediction accuracy
\end{packedenumerate}
}

